\section{The molecule rank formula (Corollary 5.7)}
\label{sec:molecule-application}

The final chapter of the molecular-conjecture program assembles the
protein-flexibility rank formula
\[
  r(G^2) \;=\; 3|V| - 6 - \operatorname{def}(\tilde G)
\]
for a graph $G$ of minimum degree at least two, where $r$ is the rank
function of the $3$-dimensional generic bar-joint rigidity matroid
(\cref{def:genericRank}) and $\tilde G = (D-1)\cdot G$ with $D = 6$. This
is Corollary~5.7 of Katoh--Tanigawa~\cite{katohTanigawa2011}; the rank
formula itself is due to Jackson--Jord\'an~\cite{jacksonJordan2008}, who
proved it equivalent to the Molecular Conjecture
(\cref{thm:molecular-conjecture}), which Katoh--Tanigawa's Theorem~5.6
(\cref{thm:theorem-55-6}) resolves.

The whole argument is arithmetic on top of the two chains already built.
No new combinatorics enters: the deficiency induction and independent-cover
bounds that Jackson--Jord\'an use to pass from the realizability
equivalence to the rank formula are replaced here by
Theorem~5.6's \emph{rank} statement, which is strictly stronger than the
$\operatorname{def} = 0$ equivalence. Two inequalities meet:

\begin{itemize}
  \item \textbf{Attainment ($\ge$).} Theorem~5.6, in its general-position
        form (\cref{thm:panel-hinge-iff-molecular}), realizes the molecular
        model of $G$ at motion-space dimension
        $D + \operatorname{def}(\tilde G)$ on centres in general position
        up to order four; the square-graph dictionary
        (\cref{thm:molecular-iff-square-bar-joint}) turns this into a
        bar-joint placement $c$ of $G^2$ with
        $\operatorname{rank} R(G^2, c) = 3|V| - 6 - \operatorname{def}(\tilde G)$,
        and the generic rank dominates the rank at any single placement.
  \item \textbf{Upper bound ($\le$).} At a placement simultaneously generic
        for row independence (\cref{def:generic-placement}) and in general
        position (\cref{lem:exists-generic-general-position}) the generic
        rank \emph{equals} the realized rank
        (\cref{lem:genericRank-eq-finrank-span}); the dictionary run in
        reverse identifies $\dim \ker R(G^2, p)$ with the molecular
        motion-space dimension, and the genericity-free lower bound
        $D + \operatorname{def}(\tilde G) \le \dim Z$
        (\cref{lem:trivial-motions-rank-bound}) caps the rank.
\end{itemize}

All the geometry is specific to $\R^3$, inherited from
\cref{sec:molecule-modelling} and \cref{sec:bar-joint-3d}.

\subsection{The generic rank dominates every placement}
\label{sec:molecule-application-domination}

\begin{lemma}[The generic rank dominates the rank at any placement]
  \label{lem:square-rank-le-genericRank}
  \lean{SimpleGraph.finrank_span_rigidityRow_le_genericRank}
  \leanok
  \uses{def:genericRank, lem:genericRigidityMatroid-indep-iff}
  Let $H$ be a graph on a finite vertex set $V$ and $p : V \to \R^d$ any
  placement. Then
  \[
    \dim \operatorname{span}
      \{\,\mathrm{rigidityRow}\,K_V\,p\,e \mid e \in E(H)\,\}
    \;\le\; r_d(H),
  \]
  i.e.\ the rank realized at $p$ never exceeds the generic rank.
  Equivalently $\operatorname{rank} R(H, p) \le r_d(H)$ for the bar-joint
  rigidity matrix at $p$.
\end{lemma}
\begin{proof}
  \leanok
  A maximal subset $I \subseteq E(H)$ on which the rigidity rows at $p$ are
  linearly independent indexes a basis of the row span, so $|I|$ is the
  dimension in question. Being row-independent at the specific placement
  $p$, $I$ is row-independent at \emph{some} placement, hence independent in
  the generic rigidity matroid $\mathcal{R}_d(V)$
  (\cref{lem:genericRigidityMatroid-indep-iff}); as $I \subseteq E(H)$ its
  cardinality is at most the rank $r_d(H) = \operatorname{rk} E(H)$.
\end{proof}

\subsection{The shadowing multigraph carrier}
\label{sec:molecule-application-carrier}

The molecular and panel-hinge machinery is stated over a multigraph with a
separate edge-label set, while the square $G^2$ and its bar-joint matroid
live over a simple graph. The two are reconciled by a multigraph that has
one edge for each adjacency of $G$ and enough spare labels for Theorem~5.6.

\begin{lemma}[A simple graph has a shadowing multigraph carrier]
  \label{lem:molecule-graph-carrier}
  \lean{SimpleGraph.shadowGraph, SimpleGraph.shadowGraph_simple,
    SimpleGraph.shadowGraph_vertexSet, SimpleGraph.shadowGraph_isLink_iff,
    SimpleGraph.card_shadowLabel_lt}
  \leanok
  \uses{def:hinge-concurrent}
  Let $G$ be a simple graph on a finite vertex set $V$. There is a
  multigraph $G'$ on $V$, spanning and simple, with an edge between
  distinct $u, v$ if and only if $u \sim_G v$, and with strictly more than
  $6(|V| - 1)$ edge labels available.
\end{lemma}
\begin{proof}
  \leanok
  Take the label type $\mathrm{Sym2}(V) \oplus \mathrm{Fin}(6(|V|-1)+1)$: the
  $\mathrm{Sym2}(V)$ summand tags each edge by its unordered endpoint pair
  (so the multigraph is simple and spanning, with an edge between distinct
  $u, v$ exactly when $u \sim_G v$), and the $\mathrm{Fin}(6(|V|-1)+1)$
  summand is pure padding, never linked to anything, whose $6(|V|-1)+1$
  labels alone already exceed the required bound --- regardless of how
  sparse $G$ is (the $\mathrm{Sym2}(V)$ summand alone can fall short, e.g.\
  for $|V| \le 10$).
\end{proof}

\begin{remark}
  Because any such carrier $G'$ is simple with the same adjacency relation
  as $G$, the crossing-edge and partition counts entering
  $\operatorname{def}$ (\cref{def:D-deficiency}) read only that relation,
  so any two such carriers have equal deficiency: $\operatorname{def}
  (\tilde{G'})$ is thereby a well-defined invariant $\operatorname{def}
  (\tilde G)$ of the simple graph $G$ alone. Not yet formalized --- the
  rank-formula assembly below fixes a single carrier throughout, so this
  invariance is not on its critical path.
\end{remark}

\subsection{The two bounds and the formula}
\label{sec:molecule-application-formula}

\begin{lemma}[Lower bound: attainment leg]
  \label{lem:molecule-rank-lower-bound}
  \uses{thm:panel-hinge-iff-molecular, thm:molecular-iff-square-bar-joint,
    lem:square-rank-le-genericRank, lem:molecule-graph-carrier}
  Let $G$ be a simple graph of minimum degree at least two on a finite
  vertex set $V$. Then $3|V| - 6 - \operatorname{def}(\tilde G) \le r(G^2)$.
\end{lemma}
\begin{proof}
  \uses{def:square-graph}
  Fix the shadowing carrier of \cref{lem:molecule-graph-carrier}. Its label
  supply meets the hypothesis of \cref{thm:panel-hinge-iff-molecular},
  which produces centres $c : V \to \R^3$ in general position up to order
  four with the molecular framework attaining motion-space dimension
  $D + \operatorname{def}(\tilde G)$. By the square-graph dictionary
  (\cref{thm:molecular-iff-square-bar-joint}) the bar-joint rigidity matrix
  $R(G^2, c)$ has $\dim \ker R(G^2, c) = D + \operatorname{def}(\tilde G)$,
  so its rank is $3|V| - 6 - \operatorname{def}(\tilde G)$. The generic rank
  dominates this single-placement rank
  (\cref{lem:square-rank-le-genericRank}).
\end{proof}

\begin{lemma}[Upper bound]
  \label{lem:molecule-rank-upper-bound}
  \uses{lem:exists-generic-general-position, lem:genericRank-eq-finrank-span,
    thm:molecular-iff-square-bar-joint, lem:trivial-motions-rank-bound,
    lem:molecule-graph-carrier}
  Let $G$ be a simple graph of minimum degree at least two on a finite
  vertex set $V$. Then $r(G^2) \le 3|V| - 6 - \operatorname{def}(\tilde G)$.
\end{lemma}
\begin{proof}
  \uses{def:square-graph, def:hinge-concurrent}
  Choose a placement $p : V \to \R^3$ that is simultaneously generic for row
  independence and in general position up to order four
  (\cref{lem:exists-generic-general-position}). At $p$ the generic rank is
  the realized rank of $G^2$ (\cref{lem:genericRank-eq-finrank-span}), so
  $r(G^2) = \operatorname{rank} R(G^2, p) = 3|V| - \dim \ker R(G^2, p)$.
  The dictionary (\cref{thm:molecular-iff-square-bar-joint}) identifies
  $\dim \ker R(G^2, p)$ with the motion-space dimension $\dim Z$ of the
  molecular framework of $G$ on the centres $p$, whose hinges are genuine
  because general position makes the centres distinct. The genericity-free
  lower bound $D + \operatorname{def}(\tilde G) \le \dim Z$
  (\cref{lem:trivial-motions-rank-bound}) then gives
  $6 + \operatorname{def}(\tilde G) \le 3|V| - r(G^2)$.
\end{proof}

\begin{corollary}[The molecule rank formula; Jackson--Jord\'an,
  Katoh--Tanigawa Corollary~5.7]
  \label{cor:molecule-rank-formula}
  \uses{lem:molecule-rank-lower-bound, lem:molecule-rank-upper-bound,
    def:genericRank, def:D-deficiency}
  Let $G = (V, E)$ be a graph of minimum degree at least two. Then
  \[
    r(G^2) \;=\; 3|V| - 6 - \operatorname{def}(\tilde G),
  \]
  where $r$ is the rank function of the $3$-dimensional generic bar-joint
  rigidity matroid (\cref{def:genericRank}) and $\tilde G = (D-1)\cdot G$
  with $D = 6$. In particular $G^2$ is generically rigid in $\R^3$ if and
  only if $\operatorname{def}(\tilde G) = 0$, i.e.\ $\tilde G$ contains six
  edge-disjoint spanning trees --- the combinatorial flexibility count for
  a molecule modelled on $G$.
\end{corollary}
\begin{proof}
  The two inequalities \cref{lem:molecule-rank-lower-bound} and
  \cref{lem:molecule-rank-upper-bound} are complementary. The rigidity
  characterization is the case $\operatorname{def}(\tilde G) = 0$, at which
  $r(G^2) = 3|V| - 6$ is the full generic rank; the tree-packing reading of
  $\operatorname{def}(\tilde G) = 0$ is the deficiency-zero specialization
  (\cref{thm:def-eq-corank}).
\end{proof}
